\documentclass[11pt]{article}

\usepackage{amsmath}
\usepackage{amsthm}
\usepackage{amsfonts}
\usepackage{amssymb}
\usepackage{booktabs}
\usepackage{hyperref}
\usepackage{xcolor}

\setlength{\textwidth}     {16.0cm}
\setlength{\textheight}    {21.0cm}
\setlength{\evensidemargin}{ 0.0cm}
\setlength{\oddsidemargin} { 0.0cm}
\setlength{\topmargin}     {-0.5cm}

\newcommand{\xkjtrial}{x^{k,j}_{\mathrm{trial}}}
\newcommand{\minimize}{\mathop{\text{Minimize}}}
\newcommand{\flow}{f_{\mathrm{low}}}
\newcommand{\dkj}{d^{k,j}}

\newtheorem{theorem}{Theorem}[section]
\newtheorem{corollary}[theorem]{Corollary}
\newtheorem{lemma}[theorem]{Lemma}
\newtheorem{definition}[theorem]{Definition}
\newtheorem{remark}[theorem]{Remark}
\newtheorem{assumption}{Assumption}
\renewcommand\theassumption{A\arabic{assumption}}

\newcounter{alg}[section]
\renewcommand{\thealg}{\thesection.\arabic{alg}}

\usepackage[mathlines]{lineno}
\renewcommand\linenumberfont{\normalfont\small\color{gray}}
\linenumbers

\begin{document}

\title{Cubic regularization for box-constrained order-value optimization\\
\large A draft of the modified theory}
\date{\today}
\maketitle

\begin{abstract}
This note develops the complexity theory for the variant of the regularized
Newton-type method for order-value optimization in which the Tikhonov term
$\tfrac{\sigma}{2}\|x - \bar{x}\|^2$ of the model is replaced by the cubic term
$\tfrac{\sigma}{3}\|x - \bar{x}\|^3$. The optimality condition
$C_{\delta,\epsilon}$ and the outer structure of the method are unchanged, so
the resulting bounds are directly comparable with those of the quadratically
regularized method. Three things change, and all three are visible in the
numerical experiments. First, the iteration complexity improves from
$O(\delta^{-2}\epsilon^{-2})$ to $O(\delta^{-3}\epsilon^{-3/2})$, so that the
cubic bound is the smaller of the two precisely when $\epsilon < \delta^{2}$.
Second, the improved order requires the curvature matrices to \emph{dominate}
the component Hessians, $B_i \succeq \nabla^2 f_i$, which makes the curvature an
essential ingredient rather than an optional one: the choice $B_i = 0$ is no
longer covered. Third, the improvement is an improvement in outer iterations
only; the number of function evaluations carries an extra
$\log_\gamma \bar{\sigma}$ factor that grows with the regularization threshold,
which is larger in the cubic case.
\end{abstract}

\section{What changes and what does not}

We keep the notation of the main paper. Given $f_1,\dots,f_m$ and
$p \in \{1,\dots,m\}$, the order-value function $f$ returns the $p$th smallest
value of $f_1(x),\dots,f_m(x)$, the problem is
$\minimize_{x \in \Omega} f(x)$ over the box
$\Omega = \{x \mid \ell \leq x \leq u\}$, and for $\delta > 0$
\[
I_\delta(x) := \{ i \mid f(x) - \delta \leq f_i(x) \leq f(x) + \delta \}
\]
is the active index set. Given $M > 0$, a family
$\mathcal{B} = \{B_i\}_{i=1}^m$ of symmetric matrices with $B_i \succeq 0$ and
$\|B_i\| \leq M$ is called \emph{admissible}.

The \emph{cubically regularized model} at $\bar{x} \in \Omega$ is
\begin{equation} \label{c-model}
\Psi_c(x; \bar{x}, \mathcal{B}, \delta, \sigma) =
\max_{i \in I_\delta(\bar{x})}
\left\lbrace \nabla f_i(\bar{x})^T(x - \bar{x})
+ \tfrac{1}{2}(x - \bar{x})^T B_i (x - \bar{x}) \right\rbrace
+ \frac{\sigma}{3}\|x - \bar{x}\|^3 .
\end{equation}
Since the regularization does not depend on $i$, writing it inside or outside
the maximum is the same function; the implementation puts it inside each
component, which is what the constraint-wise formulation of the subproblem
requires.

\begin{remark} \label{rem:strict}
$\Psi_c$ is convex, being a maximum of convex quadratics plus a convex term.
It is \emph{strictly} convex, because $x \mapsto \|x - \bar{x}\|^3$ is: if
$d_1 \neq d_2$ then
$\|\tfrac{d_1+d_2}{2}\| \leq \tfrac{1}{2}(\|d_1\| + \|d_2\|)$ and $t \mapsto
t^3$ is strictly convex on $[0,\infty)$, so equality forces
$\|d_1\| = \|d_2\|$ together with equality in the triangle inequality, hence
$d_1 = d_2$. Therefore the subproblem
\begin{equation} \label{c-subpro}
\minimize_{x \in \Omega}\, \Psi_c(x; x^k, \mathcal{B}_{k,j}, \delta, \sigma_{k,j})
\end{equation}
still has a unique global minimizer and every stationary point is global. Note
that $\Psi_c$ is \emph{not} strongly convex --- its curvature degenerates at
$x = \bar{x}$ --- so the uniqueness argument of the quadratic case, which
invoked strong convexity, has to be replaced by the above.
\end{remark}

\medskip
\noindent\refstepcounter{alg}\label{alg:cubic}\textbf{Algorithm~\thealg\ (cubic
variant).} Let $\delta > 0$, $\sigma_{\min} > 0$, $M > 0$, $\alpha \in (0,1)$,
$\gamma > 1$, and $x^0 \in \Omega$ be given. Initialize $k \leftarrow 0$.
\begin{description}
\item[Step~1.] Initialize $j \leftarrow 0$ and $\sigma_{k,j} = \sigma_{\min}$.
\item[Step~2.] Choose a curvature family $\mathcal{B}_{k,j}$ satisfying
Assumption~\ref{a6} below and compute $\xkjtrial$ as the global minimizer of
\eqref{c-subpro}.
\item[Step~3.] If
    \begin{equation} \label{c-armijo}
    f(\xkjtrial) \leq f(x^k) - \alpha \|\xkjtrial - x^k\|^3
    \end{equation}
does not hold, set $\sigma_{k,j+1} = \gamma\sigma_{k,j}$, $j \leftarrow j+1$,
and go to Step~2.
\item[Step~4.] Set $x^{k+1} = \xkjtrial$, $\sigma_k = \sigma_{k,j}$, $j_k = j$,
$k \leftarrow k+1$, and go to Step~1.
\end{description}

Throughout we write $\dkj := \xkjtrial - x^k$. Two changes with respect to the
quadratic method deserve emphasis. The acceptance test \eqref{c-armijo} is
\emph{cubic} in the step: this is what matches the order of the model error
under a Lipschitz Hessian, and it is what makes Theorem~\ref{teo:c-descent}
below work. And Step~2 no longer accepts an arbitrary admissible family: the
curvature must dominate the component Hessians, in the sense of
Assumption~\ref{a6}.

Everything else is untouched. In particular the following results, which
concern only the sign structure of the model, carry over verbatim from the main
paper.

\begin{theorem} \label{teo:c-conv}
Suppose that $f_1,\dots,f_m \in \mathcal{C}^1(\Omega)$, let $\mathcal{B}$ be
admissible and $\sigma, \delta \geq 0$. If $x^*$ is a local minimizer of
$\minimize_{x\in\Omega} f(x)$, then $x^*$ is a local minimizer of
$\minimize_{x \in \Omega} \Psi_c(x; x^*, \mathcal{B}, \delta, \sigma)$.
\end{theorem}

\begin{proof}
Identical to the chain of Theorems~3.1--3.3 of the main paper. The only place
where the shape of the regularization enters is the step in which one argues
that, if some $y^*$ makes every component of the model negative, then already
the linear parts are negative: there one uses that
$\tfrac{1}{2}(y^*-x^*)^TB_i(y^*-x^*) \geq 0$ because $B_i \succeq 0$ and that
the regularization term is nonnegative. Since
$\tfrac{\sigma}{3}\|y^*-x^*\|^3 \geq 0$ for $\sigma \geq 0$, the argument is
unchanged.
\end{proof}

Consequently the exact optimality condition of Corollary~3.4 and the
approximate condition $C_{\delta,\epsilon}$ of Definition~3.5 are
\emph{literally the same} for both methods: there exist
$\mu \in \Sigma_{|I_\delta(x)|}$ and $\nu^\ell,\nu^u \in \mathbb{R}^n_+$ with
$(\ell_i - x_i)\nu^\ell_i = (x_i - u_i)\nu^u_i = 0$ and
\begin{equation} \label{c-optcond}
\Big\| \sum_{i \in I_\delta(x)} \mu_i \nabla f_i(x)
+ \sum_{i=1}^n (\nu^u_i - \nu^\ell_i)e_i \Big\| \leq \epsilon .
\end{equation}
This is what makes the two complexity bounds comparable: they count iterations
until the same quantity falls below the same threshold.

\section{The step equation and the step bound}

Assumptions~\ref{a1}--\ref{a3} are those of the main paper.

\begin{assumption} \label{a1}
$f_1,\dots,f_m \in \mathcal{C}^1(\Omega)$.
\end{assumption}

\begin{assumption} \label{a2}
The multipliers of Corollary~\ref{cor:c-step} satisfy
$\max_i |\nu^\ell_i - \nu^u_i| \leq c_\nu$ uniformly in $k$ and $j$.
\end{assumption}

\begin{assumption} \label{a3}
$\|\nabla f_i(x)\| \leq c_\nabla$ for all $i$ and all $x \in \Omega$.
\end{assumption}

\begin{corollary} \label{cor:c-step}
Under Assumption~\ref{a1}, for every $k$ and $j = 0,\dots,j_k$ there exist
$\mu \in \Sigma_{|I_\delta(x^k)|}$ and $\nu^\ell,\nu^u \in \mathbb{R}^n_+$ with
$(\ell_i - [\xkjtrial]_i)\nu^\ell_i = ([\xkjtrial]_i - u_i)\nu^u_i = 0$ and
\begin{equation} \label{c-eqcoro}
\big( \bar{B}_{k,j} + \sigma_{k,j}\|\dkj\|\, I \big)\dkj
= \sum_{i=1}^n (\nu^\ell_i - \nu^u_i)e_i
- \sum_{i \in I_\delta(x^k)} \mu_i \nabla f_i(x^k),
\end{equation}
where $\bar{B}_{k,j} := \sum_{i \in I_\delta(x^k)} \mu_i B_{k,j,i}$ is
symmetric positive semidefinite with $\|\bar{B}_{k,j}\| \leq M$.
\end{corollary}

\begin{proof}
As in Corollary~3.6 of the main paper, the KKT conditions for the minimization
of a maximum of differentiable functions over a box give
\[
\sum_{i \in I_\delta(x^k)} \mu_i
\big[ \nabla f_i(x^k) + B_{k,j,i}\dkj \big]
+ \sigma_{k,j}\|\dkj\|\,\dkj
+ \sum_{i=1}^n (\nu^u_i - \nu^\ell_i)e_i = 0 .
\]
The function $d \mapsto \tfrac{\sigma}{3}\|d\|^3$ is continuously
differentiable on all of $\mathbb{R}^n$ with gradient $\sigma\|d\|d$, which
vanishes at $d = 0$; no nonsmoothness is introduced at the reference point.
Since $\sum_i \mu_i = 1$, the curvature terms collect into
$\bar{B}_{k,j}\dkj$, and $\bar{B}_{k,j}$, a convex combination of symmetric
positive semidefinite matrices of norm at most $M$, has the stated properties.
\end{proof}

The difference with the quadratic case is now visible: the matrix acting on the
step is $\bar{B}_{k,j} + \sigma_{k,j}\|\dkj\| I$ rather than
$\bar{B}_{k,j} + \sigma_{k,j} I$. Its smallest eigenvalue degenerates with the
step, and the resulting step bound is correspondingly weaker.

\begin{lemma} \label{lem:c-step}
Under Assumptions~\ref{a1}--\ref{a3},
\begin{equation} \label{c-ccotasub}
\|\dkj\| \leq \sqrt{\frac{c_x}{\sigma_{k,j}}}, \qquad
c_x := n\,c_\nu + c_\nabla .
\end{equation}
\end{lemma}

\begin{proof}
The right-hand side of \eqref{c-eqcoro} has norm at most
$\sum_{i=1}^n |\nu^\ell_i - \nu^u_i| + \sum_i \mu_i\|\nabla f_i(x^k)\|
\leq n c_\nu + c_\nabla = c_x$. For the left-hand side, if $\dkj \neq 0$ then
by Cauchy--Schwarz and $\bar{B}_{k,j} \succeq 0$,
\begin{align*}
\big\|(\bar{B}_{k,j} + \sigma_{k,j}\|\dkj\| I)\dkj\big\|\,\|\dkj\|
&\geq \big\langle (\bar{B}_{k,j} + \sigma_{k,j}\|\dkj\| I)\dkj, \dkj
      \big\rangle \\
&= (\dkj)^T\bar{B}_{k,j}\dkj + \sigma_{k,j}\|\dkj\|^3
\ \geq\ \sigma_{k,j}\|\dkj\|^3 ,
\end{align*}
so that
$\|(\bar{B}_{k,j} + \sigma_{k,j}\|\dkj\| I)\dkj\| \geq \sigma_{k,j}\|\dkj\|^2$.
Combining, $\sigma_{k,j}\|\dkj\|^2 \leq c_x$, which is \eqref{c-ccotasub}; the
case $\dkj = 0$ is trivial.
\end{proof}

\begin{remark}
In the quadratic method the corresponding bound is $\|\dkj\| \leq
c_x/\sigma_{k,j}$. The cubic step therefore shrinks only like
$\sigma_{k,j}^{-1/2}$. Since the mechanism that keeps the trial point inside the
$\delta$-band is precisely the smallness of the step, this is the reason why
the regularization threshold below deteriorates from $\delta^{-1}$ to
$\delta^{-2}$.
\end{remark}

\section{Sufficient descent}

The quadratic theory needed only a Lipschitz gradient. The cubic term is of
higher order than the quadratic model error $\tfrac{L}{2}\|d\|^2$ and cannot
absorb it for small steps, so a Lipschitz gradient is no longer enough: we need
the model to be a genuine second-order model, with a Lipschitz Hessian and with
curvature matrices that dominate the true ones.

\begin{assumption} \label{a4}
There is $L > 0$ with
$\|\nabla f_i(y) - \nabla f_i(x)\| \leq L\|y-x\|$ for all $i$ and
$x,y \in \Omega$.
\end{assumption}

\begin{assumption} \label{a5}
$f_1,\dots,f_m \in \mathcal{C}^2(\Omega)$ and there is $L_H > 0$ with
$\|\nabla^2 f_i(y) - \nabla^2 f_i(x)\| \leq L_H\|y-x\|$ for all $i$ and
$x,y \in \Omega$.
\end{assumption}

\begin{assumption}[dominating curvature] \label{a6}
For all $k$, all $j = 0,\dots,j_k$ and all $i \in I_\delta(x^k)$,
\[
B_{k,j,i} \succeq \nabla^2 f_i(x^k), \qquad
B_{k,j,i} \succeq 0, \qquad \|B_{k,j,i}\| \leq M .
\]
\end{assumption}

Assumption~\ref{a5} gives the standard cubic Taylor bound
\begin{equation} \label{c-taylor}
f_i(y) \leq f_i(x) + \nabla f_i(x)^T(y-x)
+ \tfrac{1}{2}(y-x)^T\nabla^2 f_i(x)(y-x)
+ \tfrac{L_H}{6}\|y-x\|^3 .
\end{equation}

\begin{remark}[on Assumption~\ref{a6}] \label{rem:a6}
The shifted exact Hessian
\[
B_{k,j,i} = \nabla^2 f_i(x^k) + \tau_{k,i} I, \qquad
\tau_{k,i} := \max\{0,\, -\lambda_{\min}(\nabla^2 f_i(x^k))\},
\]
satisfies Assumption~\ref{a6} with
$M := \max_{i}\sup_{x \in \Omega} \big( \|\nabla^2 f_i(x)\|
+ \max\{0, -\lambda_{\min}(\nabla^2 f_i(x))\} \big)$, which is finite because
$\Omega$ is compact and $f_i \in \mathcal{C}^2$. This is exactly the matrix
built by the implementation.

Two consequences must be stated plainly. First, $M$ is no longer a free
parameter that the method may impose: it is a quantity determined by the
problem. Rescaling $B_{k,j,i}$ by $M/\lambda_{\max}$ to force a prescribed
bound, as the current implementation does, destroys the domination
$B_{k,j,i} \succeq \nabla^2 f_i(x^k)$ whenever the rescaling is active, and
with it Theorem~\ref{teo:c-descent}. Second, and more importantly, the choice
$B_{k,j,i} = 0$ is admissible only if $\nabla^2 f_i(x^k) \preceq 0$. The
first-order method is therefore \emph{not} a special case of the theory of this
note: with cubic regularization, curvature is not an optional refinement but a
requirement of the analysis. Quasi-Newton families such as BFGS or SR1 are
admissible only insofar as they can be certified to dominate the component
Hessians, which the plain updates do not guarantee.
\end{remark}

\begin{theorem}[sufficient descent] \label{teo:c-descent}
Suppose that Assumptions~\ref{a1}, \ref{a3}, \ref{a4}, \ref{a5} and \ref{a6}
hold. Then, for all $k$ and $j = 0,\dots,j_k$, if
\begin{equation} \label{c-este}
\sigma_{k,j} \geq \Theta(\delta) :=
\max\left\{ \frac{L_H}{2} + 3\alpha,\ \frac{36\,c_x^3}{\delta^2},\
\frac{3\,L\,c_x}{\delta} \right\},
\end{equation}
then $\xkjtrial$ satisfies \eqref{c-armijo}.
\end{theorem}

\begin{proof}
Write $d := \dkj$ and $\sigma := \sigma_{k,j}$.

\emph{Step 1: descent for the active components.}
Let $i \in I_\delta(x^k)$. By \eqref{c-taylor} and then Assumption~\ref{a6},
\begin{align*}
f_i(\xkjtrial) &\leq f_i(x^k) + \nabla f_i(x^k)^Td
+ \tfrac{1}{2}d^T\nabla^2 f_i(x^k)d + \tfrac{L_H}{6}\|d\|^3 \\
&\leq f_i(x^k) + \nabla f_i(x^k)^Td + \tfrac{1}{2}d^TB_{k,j,i}d
+ \tfrac{L_H}{6}\|d\|^3 .
\end{align*}
Since $\xkjtrial$ minimizes \eqref{c-subpro} and
$\Psi_c(x^k; x^k,\mathcal{B}_{k,j},\delta,\sigma) = 0$, we have
$\Psi_c(\xkjtrial;\cdot) \leq 0$; as $\Psi_c$ bounds each of the components of
its maximum,
\[
\nabla f_i(x^k)^Td + \tfrac{1}{2}d^TB_{k,j,i}d + \tfrac{\sigma}{3}\|d\|^3
\leq 0 \qquad \text{for all } i \in I_\delta(x^k).
\]
Substituting,
\begin{equation} \label{c-cota-lem2}
f_i(\xkjtrial) \leq f_i(x^k)
- \Big( \frac{\sigma}{3} - \frac{L_H}{6} \Big)\|d\|^3
\leq f_i(x^k) - \alpha\|d\|^3, \qquad i \in I_\delta(x^k),
\end{equation}
the last inequality because \eqref{c-este} gives
$\sigma \geq \tfrac{L_H}{2} + 3\alpha$, i.e.\
$\tfrac{\sigma}{3} - \tfrac{L_H}{6} \geq \alpha$.

\emph{Step 2: the active set does not move out of the band.}
By Assumption~\ref{a4} and Cauchy--Schwarz, for every
$i \in \{1,\dots,m\}$,
\[
|f_i(\xkjtrial) - f_i(x^k)| \leq \|\nabla f_i(x^k)\|\,\|d\|
+ \tfrac{L}{2}\|d\|^2 \leq c_x\|d\| + \tfrac{L}{2}\|d\|^2 ,
\]
using $c_\nabla \leq c_x$. By Lemma~\ref{lem:c-step},
$\|d\| \leq (c_x/\sigma)^{1/2}$, hence
\[
|f_i(\xkjtrial) - f_i(x^k)| \leq \frac{c_x^{3/2}}{\sigma^{1/2}}
+ \frac{L\,c_x}{2\sigma} .
\]
By \eqref{c-este}, $\sigma \geq 36c_x^3/\delta^2$ gives
$c_x^{3/2}\sigma^{-1/2} \leq \delta/6$, and $\sigma \geq 3Lc_x/\delta$ gives
$Lc_x/(2\sigma) \leq \delta/6$. Therefore
\begin{equation} \label{c-teo7eq1}
|f_i(\xkjtrial) - f_i(x^k)| \leq \frac{\delta}{3}, \qquad i = 1,\dots,m .
\end{equation}

\emph{Step 3: transfer to the order-value function.}
From \eqref{c-teo7eq1}, any index with $f_i(x^k) < f(x^k) - \delta$ satisfies
$f_i(\xkjtrial) < f(x^k)$ and any index with $f_i(x^k) > f(x^k) + \delta$
satisfies $f_i(\xkjtrial) > f(x^k)$; in particular
$i_p(\xkjtrial) \in I_\delta(x^k)$ and the number of components below the band
is the same at $x^k$ and at $\xkjtrial$. Writing
$I_\delta(x^k) = \{i_1,\dots,i_r\}$ with
$f_{i_1}(x^k) \leq \cdots \leq f_{i_r}(x^k)$ and reordering as
$\{i'_1,\dots,i'_r\}$ so that
$f_{i'_1}(\xkjtrial) \leq \cdots \leq f_{i'_r}(\xkjtrial)$, the indices
$i_p(x^k)$ and $i_p(\xkjtrial)$ occupy the same position $q$ within the band.
Applying Lemma~2.1 of \cite{dunder1} with $a_j = f_{i_j}(x^k)$,
$b_j = f_{i_j}(\xkjtrial)$ and $\beta = \alpha\|d\|^3$ ---legitimate by
\eqref{c-cota-lem2}--- yields
$f_{i'_q}(\xkjtrial) \leq f_{i_q}(x^k) - \beta$, that is,
$f(\xkjtrial) \leq f(x^k) - \alpha\|d\|^3$, which is \eqref{c-armijo}.
\end{proof}

Note that Step~3 is where the cubic form of the acceptance test
\eqref{c-armijo} is forced: the ordering lemma transfers exactly the constant
$\beta$ obtained in \eqref{c-cota-lem2} from the components to $f$, and that
constant is $\alpha\|d\|^3$.

\section{Complexity}

\begin{theorem} \label{teo:c-complexity}
Suppose that Assumptions~\ref{a1}, \ref{a3}--\ref{a6} hold and that
$f(x) \geq \flow$ on $\Omega$. Put
$\bar{\sigma} := \gamma\,\Theta(\delta)$ with $\Theta(\delta)$ as in
\eqref{c-este}. Then $\sigma_k \leq \bar{\sigma}$ for every $k$, at most
\begin{equation} \label{c-cota-fcnt}
\big\lfloor 1 + \log_\gamma(\sigma_k/\sigma_{\min}) \big\rfloor
\end{equation}
function evaluations are performed per outer iteration, and the number of outer
iterations at which $C_{\delta,\epsilon}$ fails at $x^{k+1}$ is at most
\begin{equation} \label{c-cota-iter}
\left\lfloor \frac{f(x^0) - \flow}{\alpha}\,
\max\left\{ \left(\frac{2M}{\epsilon}\right)^{3},\
\left(\frac{2\bar{\sigma}}{\epsilon}\right)^{3/2} \right\} \right\rfloor .
\end{equation}
\end{theorem}

\begin{proof}
The bound $\sigma_k \leq \bar\sigma$ and \eqref{c-cota-fcnt} follow from
Theorem~\ref{teo:c-descent} and the update $\sigma_{k,j+1} = \gamma\sigma_{k,j}$
exactly as in the quadratic case: starting from $\sigma_{k,0} = \sigma_{\min}$,
at most $\lfloor 1 + \log_\gamma(\sigma_k/\sigma_{\min})\rfloor$ increases occur
before \eqref{c-este} holds, and one function evaluation is spent per increase.

Let $K$ be the set of outer indices at which $C_{\delta,\epsilon}$ fails at
$x^{k+1}$, fix $k \in K$, and write $d_k := x^{k+1} - x^k$,
$\bar{B}_k := \bar{B}_{k,j_k}$. By Corollary~\ref{cor:c-step} and the failure
of \eqref{c-optcond},
\[
\epsilon < \Big\| \sum_{i \in I_\delta(x^k)} \mu_i\nabla f_i(x^k)
+ \sum_{i=1}^n(\nu^u_i - \nu^\ell_i)e_i \Big\|
= \big\| (\bar{B}_k + \sigma_k\|d_k\| I)d_k \big\|
\leq M\|d_k\| + \sigma_k\|d_k\|^2 .
\]
Hence at least one of $M\|d_k\| > \epsilon/2$ and
$\sigma_k\|d_k\|^2 > \epsilon/2$ holds, so that, using
$\sigma_k \leq \bar\sigma$,
\[
\|d_k\| > \min\left\{ \frac{\epsilon}{2M},\
\sqrt{\frac{\epsilon}{2\bar{\sigma}}} \right\} =: s_\epsilon .
\]
Since $x^{k+1}$ satisfies \eqref{c-armijo},
$f(x^{k+1}) \leq f(x^k) - \alpha\|d_k\|^3 < f(x^k) - \alpha s_\epsilon^3$.
Summing over $k \in K$ and using $f \geq \flow$ gives
$f(x^0) - \flow > |K|\,\alpha\, s_\epsilon^3$, and
$s_\epsilon^{-3} = \max\{(2M/\epsilon)^3, (2\bar\sigma/\epsilon)^{3/2}\}$
yields \eqref{c-cota-iter}.
\end{proof}

The bound \eqref{c-cota-iter} has two branches, and which one is active is
decided by the size of the curvature relative to the target accuracy.

\begin{corollary} \label{cor:c-regime}
Under the hypotheses of Theorem~\ref{teo:c-complexity}, if
\begin{equation} \label{c-regime}
\epsilon \geq \frac{2M^2}{\bar{\sigma}},
\end{equation}
then the number of outer iterations at which $C_{\delta,\epsilon}$ fails at
$x^{k+1}$ is at most
\begin{equation} \label{c-cota-clean}
\left\lfloor \frac{f(x^0)-\flow}{\alpha}\,
\frac{(2\bar{\sigma})^{3/2}}{\epsilon^{3/2}} \right\rfloor
= O\!\left( \delta^{-3}\epsilon^{-3/2} \right).
\end{equation}
\end{corollary}

\begin{proof}
$(2M/\epsilon)^3 \leq (2\bar\sigma/\epsilon)^{3/2}$ is equivalent, on raising
both sides to the power $2/3$, to $(2M)^2 \leq 2\bar\sigma\epsilon$, i.e.\ to
\eqref{c-regime}; the first branch of \eqref{c-cota-iter} is then inactive. The
order statement follows from $\bar\sigma = \gamma\Theta(\delta)$ and
$\Theta(\delta) = O(\delta^{-2})$, so $\bar\sigma^{3/2} = O(\delta^{-3})$.
\end{proof}

\subsection*{Comparison with the quadratic method}

Table~\ref{tab:comparison} collects the two analyses. The optimality condition,
the outer loop and the inner $\sigma$-loop are identical, so the entries are
directly comparable.

\begin{table}[ht]
\centering
\begin{tabular}{lll}
\toprule
 & quadratic $\tfrac{\sigma}{2}\|d\|^2$ & cubic $\tfrac{\sigma}{3}\|d\|^3$ \\
\midrule
acceptance test & $f(x^k) - \alpha\|d\|^2$ & $f(x^k) - \alpha\|d\|^3$ \\
smoothness needed & $\nabla f_i$ Lipschitz ($L$)
   & $\nabla^2 f_i$ Lipschitz ($L_H$), and $\nabla f_i$ Lipschitz \\
curvature needed & $B_i \succeq 0$, $\|B_i\| \leq M$
   & $B_i \succeq \nabla^2 f_i$, $B_i \succeq 0$, $\|B_i\|\leq M$ \\
$B_i = 0$ admissible & yes & only if $\nabla^2 f_i \preceq 0$ \\
step bound & $\|d\| \leq c_x/\sigma$ & $\|d\| \leq (c_x/\sigma)^{1/2}$ \\
$\sigma$ threshold $\Theta(\delta)$ & $O(\delta^{-1})$ & $O(\delta^{-2})$ \\
outer iterations & $O(\delta^{-2}\epsilon^{-2})$
   & $O(\delta^{-3}\epsilon^{-3/2})$ \quad if $\epsilon \geq 2M^2/\bar\sigma$ \\
evaluations / iteration & $\lfloor 1 + \log_\gamma(\sigma_k/\sigma_{\min})\rfloor$
   & $\lfloor 1 + \log_\gamma(\sigma_k/\sigma_{\min})\rfloor$, larger $\sigma_k$ \\
\bottomrule
\end{tabular}
\caption{The quadratically and cubically regularized methods, measured against
the same approximate optimality condition $C_{\delta,\epsilon}$.}
\label{tab:comparison}
\end{table}

Three consequences are worth isolating.

\medskip
\noindent\textbf{(i) The cubic bound is smaller exactly when
$\epsilon < \delta^2$.} Ignoring constants, the two iteration bounds are
$\delta^{-2}\epsilon^{-2}$ and $\delta^{-3}\epsilon^{-3/2}$, whose ratio is
\[
\frac{\delta^{-3}\epsilon^{-3/2}}{\delta^{-2}\epsilon^{-2}}
= \frac{\epsilon^{1/2}}{\delta},
\]
which is less than one if and only if $\epsilon < \delta^2$. The cubic
regularization buys accuracy in $\epsilon$ and pays for it in $\delta$: the
narrower the active band, the larger the regularization needed to keep the
trial point inside it, and that cost is amplified by the exponent $3/2$. The
parameters of the numerical section, $\delta = 10^{-1}$ and
$\epsilon = 10^{-4}$, satisfy $\epsilon < \delta^2 = 10^{-2}$ comfortably,
which is the regime in which the cubic bound is the better of the two.

\medskip
\noindent\textbf{(ii) The curvature is what limits the improved regime.} In the
quadratic analysis the bound $M$ on the curvature enters only through the
constant $(M + \gamma\Theta(\delta))^2$, because $M$ and $\sigma$ multiply the
step in the same way. Under cubic regularization the residual is
$M\|d\| + \sigma\|d\|^2$, so the two terms scale differently and $M$ enters the
\emph{order}: below the threshold $\epsilon < 2M^2/\bar\sigma$ the bound
degrades to $O(\epsilon^{-3})$. The improved rate is therefore a
moderate-accuracy statement, valid on the window
\[
\frac{2M^2}{\bar{\sigma}} \leq \epsilon < \delta^2 ,
\]
which is nonempty whenever $2M^2 < \bar\sigma\delta^2$; with
$\bar\sigma \geq 36\gamma c_x^3\delta^{-2}$ this holds as soon as
$M^2 < 18\gamma c_x^3$, a mild condition. One would like to shrink $M$ to widen
the window, but Assumption~\ref{a6} forbids it: $M$ must be large enough for
the $B_i$ to dominate the component Hessians. This is a genuine tension of the
cubic variant and it should be stated rather than hidden.

\medskip
\noindent\textbf{(iii) The gain is in iterations, not in evaluations.} The
total number of function evaluations is at most the number of outer iterations
times $\lfloor 1 + \log_\gamma(\bar\sigma/\sigma_{\min})\rfloor$. The cubic
method has a larger $\bar\sigma$ ---$O(\delta^{-2})$ against
$O(\delta^{-1})$--- so it spends more inner increases of $\sigma$ per outer
iteration, and the improvement in evaluations is smaller than the improvement
in iterations by a factor of order
$\log_\gamma(\Theta_{\text{cub}}/\Theta_{\text{quad}})
= \log_\gamma O(\delta^{-1})$. This is precisely the pattern of the
experiments: on the nonlinear problem the proposed method improves the mean
iteration count by $4.44\times$ but the mean evaluation count only by
$2.76\times$, and on the problem affine in the parameters it improves
iterations by $1.22\times$ while \emph{losing} $8\%$ in evaluations.

\section{What the implementation would have to change}

Three points connect the theory above with the code that produced the numbers.
The first has been settled; the remaining two must be reconciled before the
numbers can be presented as illustrating these theorems.

\begin{enumerate}
\item \textbf{The acceptance test} \emph{(done).} Step~3 of the implementation
now uses the cubic test \eqref{c-armijo} in the second-order branch, and retains
the quadratic test $f(\xkjtrial) \leq f(x^k) - \alpha\|\xkjtrial - x^k\|^2$ in
the first-order branch, which is the published method. Both problems were
rerun after the change: the mean iteration and evaluation counts are
\emph{identical} to those obtained with the quadratic test, on every value of
$o$ and on both problems. This is what one expects, since $\alpha = 10^{-8}$ and
$\|d\| < 1$ throughout, so the two right-hand sides differ by less than the gap
between consecutive accepted and rejected trial points; but it is worth having
verified rather than argued.

\item \textbf{The rescaling of the curvature.} The implementation shifts the
exact Hessian to positive definiteness and then rescales it by
$M/\lambda_{\max}$ when $\lambda_{\max} > M$, with $M = 1$. The shift is what
Assumption~\ref{a6} asks for; the rescaling breaks it. Either the rescaling is
removed and $M$ is redefined as in Remark~\ref{rem:a6} as the bound that the
shifted Hessians happen to satisfy on $\Omega$, or Theorem~\ref{teo:c-descent}
does not apply to the runs. Removing it changes the numbers and the experiments
would have to be rerun.

\item \textbf{The first-order comparison.} The baseline in the experiments is
the quadratically regularized method with $B_i = 0$, which is covered by the
theory of the main paper but not by the theory of this note. That is the
correct comparison to report ---it is the published method--- but the two
columns of the tables are then governed by two different sets of assumptions,
and the text must say so. In particular, the cubically regularized method with
$B_i = 0$, which performed best of all in the exploratory runs, is not covered
by either analysis.
\end{enumerate}

\begin{thebibliography}{9}
\bibitem{dunder1} Order-value optimization ordering lemma; see Lemma~2.1 of the
reference cited as \texttt{dunder1} in the main paper.
\bibitem{alvarezovo} The first-order regularized method for order-value
optimization generalized by the main paper.
\end{thebibliography}

\end{document}
